\documentclass[11pt]{amsart}
\usepackage{amsmath,amssymb,amsthm}
\usepackage{enumitem}
\usepackage{booktabs,array}
\usepackage{hyperref}
\hypersetup{colorlinks=true,linkcolor=black,citecolor=black,urlcolor=blue}

\newtheorem{theorem}{Theorem}[section]
\newtheorem{proposition}[theorem]{Proposition}
\newtheorem{lemma}[theorem]{Lemma}
\newtheorem{corollary}[theorem]{Corollary}
\theoremstyle{definition}
\newtheorem{definition}[theorem]{Definition}
\theoremstyle{remark}
\newtheorem{remark}[theorem]{Remark}
\newtheorem{example}[theorem]{Example}

\newcommand{\Gp}{G}
\newcommand{\dd}{\mathrm{d}}
\newcommand{\Fp}{\mathbb{F}_p}
\newcommand{\Q}{\mathbb{Q}}
\newcommand{\Z}{\mathbb{Z}}
\newcommand{\Zp}{\mathbb{Z}_{(p)}}
\newcommand{\leg}[2]{\left(\frac{#1}{#2}\right)}
\newcommand{\fl}[1]{\left\lfloor #1 \right\rfloor}

\title[A finite Euler--Clausen transfer for Ap\'ery-like sums]{A finite Euler--Clausen transfer for Ap\'ery-like sums,\\ with an application to a conjecture of Zhi-Hong Sun}

\author{Huimin Zheng}
\address{School of Information \& Network Engineering, Anhui Science and Technology University, Fengyang, Anhui 233100, People's Republic of China}
\email{zhenghm@ahstu.edu.cn}
\thanks{ORCID: 0009-0004-8605-2713}

\subjclass[2020]{Primary 11A07; Secondary 11B65, 11E25, 33C20, 11F33}
\keywords{Ap\'ery-like numbers; supercongruence; truncated hypergeometric series; Clausen's identity; Ramanujan-type series for $1/\pi$; binary quadratic forms}

\begin{document}

\begin{abstract}
Let $G_n=\sum_{k=0}^n\binom{2k}{k}^2\binom{2n-2k}{n-k}4^{\,n-k}$ be the
Ap\'ery-like sequence of Almkvist--Zudilin. We study the finite sums
$\sum_{n<p}\binom{2n}{n}G_n/M^n$ and their first moments
$\sum_{n<p}n\binom{2n}{n}G_n/M^n$ modulo powers of an odd prime $p$.
The generating series
$H(z)=\sum_{n\ge0}\binom{2n}{n}G_nz^n/64^n$ satisfies the exact identity
$H(z)=K(4z(1-z))$, where
$K(w)={}_3F_2(\frac14,\frac12,\frac34;1,1;w)$.
We prove that truncating this identity at degree $p$ commutes with
specialization at the center $z=\frac12$ up to $p^3$, and that its
Euler derivative at any $p$-integral point with $1-z$ a $p$-unit is
correct up to $p^2$; in both cases the first defect is computed
explicitly and shown to possess a reflection symmetry, a
constant-coefficient forcing equation and a short hypergeometric kernel.
Combining these finite transfers with published special values of
truncated triple-binomial sums (Sun, Sun--Ye) and with an explicit
specialization of the general CM theorem of
Chisholm--Deines--Long--Nebe--Swisher, we prove all nine formulas of
Conjecture~2.4 of Z.-H.~Sun [\emph{New congruences involving
Ap\'ery-like numbers}, arXiv:2004.07172v2]: three pairs of weighted
congruences modulo $p^2$ with denominators
$(72,576)$, $(48,-192)$, $(63,-4032)$, and the unweighted congruence with
denominator $128$ modulo $p^3$.
Selected algebraic components of the argument have been checked in Lean.
\end{abstract}

\maketitle

\section{Introduction}

For a nonnegative integer $n$, the Ap\'ery-like sequence studied in this
paper is
\begin{equation}\label{eq:Gdef}
G_n=\sum_{k=0}^n\binom{2k}{k}^2\binom{2n-2k}{n-k}4^{\,n-k}
\end{equation}
(OEIS A143583; $1,12,164,2352,34596,\dots$). It is one of the sporadic
second-kind Ap\'ery-like sequences of Almkvist--Zudilin \cite{AZ} and
satisfies
\[
(n+1)^2G_{n+1}=(32n(n+1)+12)G_n-256n^2G_{n-1}\qquad(n\ge1).
\]
Zhi-Hong Sun~\cite{Sun24} initiated a systematic study of the
$p$-adic behavior of the sums
\begin{equation}\label{eq:sums}
\sum_{n=0}^{p-1}\frac{\binom{2n}{n}G_n}{M^n},
\qquad
\sum_{n=0}^{p-1}\frac{n\binom{2n}{n}G_n}{M^n},
\end{equation}
and stated there a number of precise conjectures. The present paper is
concerned with the following one \cite[Conjecture~2.4]{Sun24}.

\begin{quote}\small
\textbf{Conjecture (Z.-H.~Sun).} \emph{Let $p>3$ be a prime. Then}
\[
\begin{aligned}
\sum_{n=0}^{p-1}\frac{n\binom{2n}{n}G_n}{72^n}
&\equiv-8\sum_{n=0}^{p-1}\frac{n\binom{2n}{n}G_n}{576^n}\\
&\equiv\begin{cases}
4x^2-3p\pmod{p^2}&\text{if }p=x^2+4y^2\equiv1\pmod4,\\
p\pmod{p^2}&\text{if }p\equiv3\pmod4,
\end{cases}
\end{aligned}
\]
\[
\begin{aligned}
\sum_{n=0}^{p-1}\frac{n\binom{2n}{n}G_n}{48^n}
&\equiv4\sum_{n=0}^{p-1}\frac{n\binom{2n}{n}G_n}{(-192)^n}\\
&\equiv\begin{cases}
3p-4x^2\pmod{p^2}&\text{if }p=x^2+3y^2\equiv1\pmod3,\\
-p\pmod{p^2}&\text{if }p\equiv2\pmod3,
\end{cases}
\end{aligned}
\]
\[
\begin{aligned}
\sum_{n=0}^{p-1}\frac{n\binom{2n}{n}G_n}{63^n}
&\equiv64\sum_{n=0}^{p-1}\frac{n\binom{2n}{n}G_n}{(-4032)^n}\\
&\equiv\begin{cases}
-8(4x^2-3p)\pmod{p^2}&\text{if }p=x^2+7y^2\equiv1,2,4\pmod7,\\
-8p\pmod{p^2}&\text{if }p\equiv3,5,6\pmod7,
\end{cases}
\end{aligned}
\]
\emph{and}
\[
\sum_{n=0}^{p-1}\frac{\binom{2n}{n}G_n}{128^n}\equiv
\begin{cases}
4x^2-2p-\dfrac{p^2}{4x^2}\pmod{p^3}&\text{if }p=x^2+2y^2\equiv1,3\pmod8,\\[2mm]
\dfrac{p^2}{3}\binom{\lfloor p/4\rfloor}{\lfloor p/8\rfloor}^{-2}\pmod{p^3}
&\text{if }p\equiv5\pmod8,\\[2mm]
-\dfrac{3p^2}{2}\binom{\lfloor p/4\rfloor}{\lfloor p/8\rfloor}^{-2}\pmod{p^3}
&\text{if }p\equiv7\pmod8.
\end{cases}
\]
\end{quote}

(The third pair involves the denominator $63$, so at $p=7$ the sums are
not $p$-integral and the statement is understood for $p\ne7$.)

\medskip

\textbf{Why these sums are natural, and what was missing.}
Sums of the shape \eqref{eq:sums} are the Ap\'ery-side incarnation of a
classical theme going back to Van Hamme \cite{vH} and
Rodriguez-Villegas~\cite{RV}: truncated hypergeometric series at CM
points are expected to be congruent, modulo $p^2$ or $p^3$, to quantities
built from binary quadratic forms. For the \emph{triple binomial sums}
\begin{equation}\label{eq:TU}
T_p(d)=\sum_{k=0}^{p-1}\frac{a_k}{d^k},
\qquad
U_p(d)=\sum_{k=0}^{p-1}\frac{k\,a_k}{d^k},
\qquad
a_k=\binom{2k}{k}^2\binom{4k}{2k}=\frac{(4k)!}{(k!)^4},
\end{equation}
an extensive theory now exists: Sun evaluated $U_p(648)$ and $U_p(-144)$
modulo $p^2$ \cite[Theorems~2.9, 2.10]{SunBQF}, Sun--Ye evaluated
$T_p(d)$ modulo $p^3$ at several CM points including $d=256$ and
$d=-3969$ \cite[Theorems~1.5, 1.8]{SunYe}, and a general theorem of
Chisholm--Deines--Long--Nebe--Swisher \cite[Theorem~1]{CDLNS} evaluates
truncated CM Ramanujan-type sums modulo $p^2$, covering both ordinary
and supersingular primes.

The sums \eqref{eq:sums}, however, involve the Ap\'ery-like numbers
$G_n$ themselves, not the triple binomial coefficients $a_k$. What
separates the two is a \emph{finite transformation problem}: the
generating series of $\binom{2n}{n}G_n/64^n$ is related to the
${}_3F_2$ series $K(w)$ by the quadratic substitution $w=4z(1-z)$
(Theorem~\ref{thm:A} below), but \emph{truncating} an identity of formal
power series at degree $p$ and then specializing at a point is not
congruence-compatible in general: the discarded tail $k>(p-1)/2$ has
coefficients divisible only by $p^2$ (Lemma~\ref{lem:val}), and at a
generic point this residual $p^2$ is exactly what one sees. The content
of the present paper is that at the center $z=\frac12$ a further factor
$p$ appears (so the center comparison holds modulo $p^3$), while the
first $p^2$-defect at a generic point is an explicit polynomial
$R_p(z)$ whose \emph{odd reflection symmetry} is responsible for the
pairing relations in Sun's conjecture. Controlling this defect ---
not evaluating any new CM period --- is the technical heart of the
matter.

\medskip

\textbf{Main results.} Throughout, $p$ is an odd prime,
$\Zp$ denotes the localization of $\Z$ at $p$, and
\[
H_p(z)=\sum_{n=0}^{p-1}\frac{\binom{2n}{n}G_n}{64^n}z^n,
\qquad
K_p(w)=\sum_{k=0}^{p-1}\frac{a_k}{256^k}w^k .
\]

\begin{itemize}
\item \textbf{Theorem~\ref{thm:A}} (formal identity).
$H(z)=K(4z(1-z))$ in $\Q[[z]]$; equivalently $H(z)=B(z)^2$ with
$B(z)={}_2F_1(\frac14,\frac34;1;z)$.

\item \textbf{Theorem~\ref{thm:B}} (center bridge). For \emph{every} odd
prime $p$,
\[
\sum_{n=0}^{p-1}\frac{\binom{2n}{n}G_n}{128^n}
\equiv\sum_{k=0}^{p-1}\frac{a_k}{256^k}\pmod{p^3}.
\]

\item \textbf{Theorem~\ref{thm:Bp}} (first defect at the center). The
quotient of the difference by $p^3$ is given modulo $p$ by an explicit
finite double sum \eqref{eq:Delta}; in particular the center comparison
does \emph{not} lift to $p^4$ in general.

\item \textbf{Theorem~\ref{thm:C}} (Euler transfer).
$H_p(z)-K_p(4z(1-z))\in p^2\Zp[z]$, and the first defect has the form
$p^2z^pR_p(z)\pmod{p^3}$ with $R_p(1-z)=-R_p(z)$.

\item \textbf{Corollary~\ref{cor:pair}} (pairing and one-sided
evaluations). For every $z$ with $z,1-z\in\Zp^\times$,
\begin{gather*}
(1-z)\,\mathcal W_p(z)+z\,\mathcal W_p(1-z)\equiv0\pmod{p^2},\\
\mathcal W_p(z)\equiv\frac{1-2z}{1-z}\,(\theta_wK_p)(w(z))\pmod{p^2},
\end{gather*}
where $\mathcal W_p(z)=zH_p'(z)$ and $\theta_w=w\,\dd/\dd w$.

\item \textbf{Theorem~\ref{thm:D}} (structure of the defect).
$\mathcal SR_p=\varepsilon/4$ in $\Fp[z]$ for an explicit third-order
operator $\mathcal S$ and $\varepsilon=\leg{-1}{p}$; $R_p$ factors as
$(1-2z)J_p(4z(1-z))$ with an explicit hypergeometric kernel $J_p$ of
degree $\lfloor(p-3)/4\rfloor$; and $R_p$ is determined by a polynomial
integral invariant with a characteristic-$p$ correction term $z^p$.

\item \textbf{Theorems~\ref{thm:128split} and \ref{thm:128inert}}
(denominator $128$ modulo $p^3$).
The last displayed formula of the conjecture holds: the split branch by
combining Theorem~\ref{thm:B} with \cite[Theorem~1.5]{SunYe}, the inert
branches by combining it with a jet comparison
(Proposition~\ref{prop:jet}) and an exact rational identity
$\mathcal Q_a(\tfrac12)=-\bigl(2^{a-1}/(a\binom{a-1}{(a-1)/2})\bigr)^2$
for odd $a$ (Theorem~\ref{thm:E}) proved by a symmetric-square method.

\item \textbf{Theorem~\ref{thm:final}}. All nine formulas of Sun's
Conjecture~2.4 hold.
\end{itemize}

\textbf{What is new and what is cited.} The finite transfers
(Theorems~\ref{thm:B}, \ref{thm:Bp}, \ref{thm:C}, \ref{thm:D}),
the symmetric-square identity (Theorem~\ref{thm:E}) and the assembly are
proved here. All CM special values of the triple binomial sums are taken
from the literature: \cite[Theorems~2.9, 2.10 and Remark~2.3]{SunBQF},
\cite[Theorems~1.5, 1.8]{SunYe}, \cite[Theorem~5.1]{SunL2}, and the
general CM theorem \cite[Theorem~1]{CDLNS}, whose specialization to the
$65k+8$ series at $\lambda=-256/3969$ we verify in detail, including the
two primes $p=5,13$ excluded by its unit hypotheses
(Proposition~\ref{prop:ram65}). The classical identity
$\sum_{k\ge0}(65k+8)a_k/(-3969)^k=9\sqrt7/\pi$ \cite[Table~2]{Guillera}
is used only to identify the parameters of that theorem.
Mao--Liu \cite[Theorem~1.10]{MaoLiu} have obtained the \emph{quadratic}
and $1/(2n-1)$ moments of the denominator-$128$ sum modulo $p^2$ by a
different reduction; their results neither contain nor follow from the
present ones (different weights, different modulus). The unweighted
denominator-$128$ sum modulo $p^3$ and the linear moments at
$72,576,48,-192,63,-4032$ do not appear in the literature we surveyed;
we make no claim of priority beyond this.

Other recent results on the same Ap\'ery-like family concern parameter
moments without the additional central-binomial weight
\cite{MaoYang,Mao}, or Christoffel--Darboux identities and endpoint
congruences for the polynomial families $g_n(x)$ and $v_n(x)$
\cite{SunCD}.  These works provide useful context but do not furnish the
finite quadratic pullback used here.  This distinction is important:
the new step in the present argument is not another isolated CM
evaluation, but control of the error created by truncating that pullback.

\medskip

\textbf{Outline.} Section~2 proves the formal identity.
Section~3 proves the center bridge and its first defect.
Section~4 proves the Euler transfer and the pairing corollary.
Section~5 analyses the defect polynomial $R_p$.
Section~6 assembles the three weighted pairs from the cited special
values, with the exceptional primes settled directly.
Section~7 proves the denominator-$128$ formulas modulo $p^3$.
Section~8 states the final corollary and some remarks, including an
obstruction to a naive finite Pfaff transformation.
Appendix~A collects the finite exceptional-prime certificates;
Appendices~B and~C give the self-contained proofs of the jet comparison
and of the symmetric-square identity; Appendix~D describes the
computational checks.

\medskip

\textbf{Acknowledgments.} The proofs and the initial manuscript draft
were developed with substantial assistance from OpenAI's Codex. The
author reviewed the mathematical arguments, computational checks,
references, and final presentation and assumes responsibility for the
contents of the article.

\medskip
\noindent\textbf{Data availability.} No datasets were generated or
analysed in this work.  The finite computations described in
Appendix~\ref{app:D} are reproducible from the displayed definitions.

\medskip
\noindent\textbf{Competing interests.} The author declares no competing
interests.

\section{Notation and the generating-function identity}\label{sec:2}

Throughout $p$ is an odd prime and $v_p$ the $p$-adic valuation. We set
\[
a_k=\binom{2k}{k}^2\binom{4k}{2k}=\frac{(4k)!}{(k!)^4},
\qquad
c_k=\frac{a_k}{256^k},
\qquad
h_n=\frac{\binom{2n}{n}G_n}{64^n},
\]
\[
H(z)=\sum_{n\ge0}h_nz^n,\qquad
K(w)=\sum_{k\ge0}c_kw^k
={}_3F_2\Bigl(\tfrac14,\tfrac12,\tfrac34;1,1;w\Bigr),
\]
and denote by $H_p$, $K_p$ the truncations at degree $p-1$. Write
\[
s=z(1-z),\qquad t=1-2z,\qquad w=w(z)=4s .
\]

\begin{lemma}[valuation]\label{lem:val}
For $0\le k<p$ one has
$v_p(a_k)=v_p(c_k)=\fl{4k/p}$; in particular $v_p(c_k)\ge2$ for
$k>(p-1)/2$, and $v_p(c_k)\ge3$ for $k\ge(3p+1)/4$.
\end{lemma}

\begin{proof}
Since $k<p$, $v_p(k!)=0$, and $4k\le4(p-1)<p^2$ because
$p^2-4(p-1)=(p-2)^2>0$; hence $v_p(a_k)=v_p((4k)!)=\fl{4k/p}$.
\end{proof}

\begin{theorem}[formal identity]\label{thm:A}
In $\Q[[z]]$,
\[
H(z)=K(4z(1-z)).
\]
Equivalently $H(z)=B(z)^2$, where
\[
B(z)={}_2F_1\Bigl(\frac14,\frac34;1;z\Bigr)
=\sum_{n\ge0}\binom{2n}{n}\binom{4n}{2n}\frac{z^n}{64^n}.
\]
\end{theorem}

\begin{proof}
\emph{Step 1: a second-order equation for $g(x)=\sum_{n\ge0}G_nx^n$.}
Since $(1-16x)^{-1/2}=\sum_{m\ge0}\binom{2m}{m}(4x)^m$ and
${}_2F_1(\frac12,\frac12;1;16x)=\sum_{k\ge0}\binom{2k}{k}^2x^k$,
the defining convolution \eqref{eq:Gdef} is exactly the Cauchy product
\begin{equation}\label{eq:cauchy}
g(x)=(1-16x)^{-1/2}\,{}_2F_1\Bigl(\tfrac12,\tfrac12;1;16x\Bigr).
\end{equation}
The hypergeometric factor satisfies the Gauss equation
$x(1-16x)u''+(1-32x)u'-4u=0$; substituting
$u=(1-16x)^{1/2}g$ and clearing the invertible factor gives
\[
x(1-16x)^2g''+(1-48x)(1-16x)g'+(256x-12)g=0,
\]
i.e. coefficientwise
\begin{equation}\label{eq:Grec}
(n+1)^2G_{n+1}=(32n(n+1)+12)G_n-256n^2G_{n-1}\qquad(n\ge1).
\end{equation}

\emph{Step 2: a third-order equation for $H$.}
Inserting $G_n=64^nh_n/\binom{2n}{n}$ into \eqref{eq:Grec} and using
$\binom{2n+2}{n+1}/\binom{2n}{n}=2(2n+1)/(n+1)$ gives the exact
three-term recurrence
\begin{equation}\label{eq:hrec}
(n+1)^3h_{n+1}
=(2n+1)\Bigl(n(n+1)+\tfrac38\Bigr)h_n
-\frac{n(2n-1)(2n+1)}{4}h_{n-1}.
\end{equation}
Hence, with $\theta=z\,\dd/\dd z$, $H$ is annihilated by
\[
\theta^3-\frac z8(2\theta+1)\bigl(8\theta(\theta+1)+3\bigr)
+\frac{z^2}4(\theta+1)\bigl(4(\theta+1)^2-1\bigr)=z\mathcal S,
\]
where
\begin{equation}\label{eq:S}
\mathcal S=s^2\partial_z^3+3st\,\partial_z^2
+\Bigl(t^2-\tfrac{11}{4}s\Bigr)\partial_z-\tfrac38t .
\end{equation}

\emph{Step 3: symmetric square.} $B$ satisfies the Gauss equation
$\mathcal D B=0$ with
$\mathcal D=z(1-z)\partial_z^2+(1-2z)\partial_z-\frac3{16}$, and $\mathcal S$
is the corresponding symmetric-square operator: for any $f,v$ one has
the exact product identity
\begin{equation}\label{eq:symmsq}
\mathcal S(fv)=sf(\mathcal Dv)'+(3sf'+tf)\mathcal Dv
+sv(\mathcal Df)'+(3sv'+tv)\mathcal Df ,
\end{equation}
so $\mathcal S(B^2)=0$. Both $H$ and $B^2$ have constant term $1$, and
in \eqref{eq:hrec} the leading coefficient $(n+1)^3$ is invertible in
$\Q$; hence a solution of $\mathcal SY=0$ in $\Q[[z]]$ is uniquely
determined by its constant term, and $H=B^2$.

\emph{Step 4: the quadratic substitution.}
The ratio $c_k/c_{k-1}=\frac{(k-\frac14)(k-\frac12)(k-\frac34)}{k^3}$
shows that $K$ is annihilated by
\[
\mathcal M=w^2(1-w)\partial_w^3+w\Bigl(3-\tfrac92w\Bigr)\partial_w^2
+\Bigl(1-\tfrac{51}{16}w\Bigr)\partial_w-\tfrac3{32}.
\]
Writing $s=z(1-z)$, $t=1-2z$ and differentiating
$K(w(z))$ formally (using $w'=4t$, $t^2=1-w$) gives the exact
operator identity, valid without applying $\mathcal MK=0$,
\begin{equation}\label{eq:pullback}
\mathcal S_z\bigl[K(w(z))\bigr]=4t\,(\mathcal M_wK)(w(z)).
\end{equation}
Hence $\mathcal S(K\circ w)=0$. Since $K(w(z))$ has constant term $1$,
the same uniqueness argument gives $H=K\circ w$. The equivalence of the
two forms is Clausen's identity
${}_2F_1(\frac14,\frac34;1;z)^2
={}_3F_2(\frac14,\frac12,\frac34;1,1;4z(1-z))$.
\end{proof}

\begin{remark}
The identity $H=B^2$ is of classical Clausen type
(compare \cite{CTYZ}); we included the operator derivation above because
every subsequent argument truncates \eqref{eq:pullback} at degree $p$,
and the truncation error is computed from the same two operators.
\end{remark}

For later use we record the finite form of the Clausen step. Let
$b_k=\binom{2k}{k}\binom{4k}{2k}/64^k$ and
$B_p(t)=\sum_{k<p}b_kt^k$.

\begin{proposition}[finite Clausen square]\label{prop:fclausen}
For every odd prime $p$,
\[
K_p\bigl(4t(1-t)\bigr)\equiv B_p(t)^2\pmod{p^2}
\qquad\text{in }\Zp[t].
\]
\end{proposition}

\begin{proof}
Coefficients of degree $<p$ agree exactly by Theorem~\ref{thm:A}.
On the left, terms with $k>(p-1)/2$ vanish modulo $p^2$ by
Lemma~\ref{lem:val}, and $k\le(p-1)/2$ contributes degree $\le p-1$
only. On the right, the coefficient of degree $n\ge p$ is
$\sum_{i+j=n}b_ib_j$. For $0\le i<p$,
\[
v_p(b_i)=v_p\Bigl(\tfrac{(4i)!}{(2i)!}\Bigr)=\fl{\tfrac{4i}{p}}-\fl{\tfrac{2i}{p}}
=\begin{cases}0,&i<p/4,\\1,&p/4<i<3p/4,\\2,&3p/4<i<p.\end{cases}
\]
If $i+j\ge p$ then either one index is $<p/4$ and the other $>3p/4$, or
both indices are $\ge p/4$; in both cases $v_p(b_ib_j)\ge2$.
\end{proof}

\section{The center bridge modulo \texorpdfstring{$p^3$}{p-cubed}}\label{sec:3}

\begin{theorem}[center bridge]\label{thm:B}
For every odd prime $p$,
\[
\sum_{n=0}^{p-1}\frac{\binom{2n}{n}G_n}{128^n}
\equiv\sum_{k=0}^{p-1}\frac{a_k}{256^k}\pmod{p^3}.
\]
\end{theorem}

\begin{proof}
Truncate the identity of Theorem~\ref{thm:A} at degree $<p$ and then set
$z=\frac12$. Since $w(z)=4z(1-z)$ has vanishing order $1$, the terms
$k\ge p$ of $K$ do not meet the coefficients of degree $<p$, so
\begin{equation}\label{eq:Rpk}
H_p\Bigl(\tfrac12\Bigr)=\sum_{k=0}^{p-1}c_kR_{p,k},
\qquad
R_{p,k}=\sum_{j=0}^{\min(k,p-1-k)}(-1)^j\binom{k}{j}2^{\,k-j}\in\Z .
\end{equation}
Everything takes place in a finite polynomial ring; no infinite series is
evaluated at the degenerate point $w=1$.

If $2k<p$ the full binomial expansion of $(2-1)^k$ is present and
$R_{p,k}=1$ exactly.

If $2k>p$, write $r=p-1-k$, so $0\le r<k<p$. For $0\le j\le r$ the
factor $j!$ is a $p$-unit and $k\equiv-1-r\pmod p$, hence
\[
\binom{k}{j}\equiv(-1)^j\binom{r+j}{j}\pmod p .
\]
Combining with the weight $2^{k-j}$ and the finite identity
\begin{equation}\label{eq:pascal}
\sum_{j=0}^{r}\binom{r+j}{j}2^{-j}=2^r
\qquad\text{in }\Q
\end{equation}
(proof: if $S_r$ denotes the left side, Pascal's rule gives
$S_r=S_{r-1}+2^{-r}\binom{2r-1}{r}+\frac12\bigl(S_r-2^{-r}\binom{2r}{r}\bigr)$,
and $\binom{2r}{r}=2\binom{2r-1}{r}$ forces $S_r=2S_{r-1}$), we obtain
\[
R_{p,k}\equiv2^{k+r}=2^{p-1}\equiv1\pmod p .
\]
By Lemma~\ref{lem:val}, the high range $2k>p$ has $v_p(c_k)\ge2$, so each
summand $c_k(R_{p,k}-1)$ lies in $p^3\Zp$; the low-range summands vanish
exactly. Summing over $k$ gives the claim, for all odd $p$ including
$p=3$.
\end{proof}

\begin{remark}
The extra factor $p$ at the center has a visible combinatorial origin:
the quadratic coordinate $4z(1-z)$, the reflection
$k\leftrightarrow p-1-k$ of the truncated binomial expansion, and the
carry valuation $\fl{4k/p}$ cooperate. No such extra factor exists at a
generic point (Section~\ref{sec:4}).
\end{remark}

The theorem is sharp: the difference is exactly divisible by $p^3$ in
the smallest cases, and its quotient modulo $p$ is as follows.

\begin{theorem}[first defect at the center]\label{thm:Bp}
Let $p$ be an odd prime, $m=(p-1)/2$, and
$\Delta_p=p^{-3}\bigl(H_p(\frac12)-K_p(1)\bigr)\in\Zp$. Then
\begin{equation}\label{eq:Delta}
\Delta_p\equiv
\sum_{k=(p+1)/2}^{\lfloor(3p-1)/4\rfloor}
\frac{a_k}{p^2}\,256^{-k}
\sum_{j=p-k}^{m}\frac{\binom{2j}{j}}{2j\,4^j}
\pmod p .
\end{equation}
For $p=3,5,7,11,13,17$ the values are $1,1,2,1,5,12$ respectively; in
particular the center comparison does not lift to $p^4$ in general.
\end{theorem}

\begin{proof}
For $m<k<p$ write $r=p-1-k$. Since $r+i\le2r<p$, all denominators in
\[
(-1)^j\binom{k}{j}
=\binom{r+j}{j}\prod_{i=1}^{j}\Bigl(1-\frac{p}{r+i}\Bigr)
\]
are $p$-units, and expanding to first order,
\[
R_{p,k}\equiv2^k\bigl(2^r-pU_r\bigr)\pmod{p^2},
\qquad
U_r=\sum_{j=0}^{r}\binom{r+j}{j}2^{-j}(H_{r+j}-H_r),
\]
with $H_n$ the harmonic numbers. The polynomial
$P_r(X)=\sum_{j=0}^r\binom{r+X+j}{j}2^{-j}$ satisfies
$P_r(X)=2P_{r-1}(X)+2^{-r}\frac{X}{r+X}\binom{2r+X-1}{r}$
by Pascal's rule; differentiating at $X=0$ (a finite polynomial
operation) gives
$U_r=2U_{r-1}+\binom{2r}{r}/(r2^{r+1})$, hence the exact identity
\[
2^{-r}U_r=L_r:=\sum_{j=1}^{r}\frac{\binom{2j}{j}}{2j\,4^j}.
\]
Therefore, with $q_p(2)=(2^{p-1}-1)/p$,
\begin{equation}\label{eq:Rexp}
\frac{R_{p,k}-1}{p}\equiv q_p(2)-L_r\pmod p .
\end{equation}
At the boundary $k=r=m$ one has $R_{p,m}=1$ exactly, whence
$q_p(2)\equiv L_m\pmod p$ and the bracket in \eqref{eq:Rexp} is
$L_m-L_r=\sum_{j=r+1}^{m}\binom{2j}{j}/(2j4^j)$. Finally
$\Delta_p=\sum_{k=m+1}^{p-1}\frac{a_k}{p^2}256^{-k}\frac{R_{p,k}-1}{p}$;
by Lemma~\ref{lem:val}, $\frac{a_k}{p^2}$ is integral, and it is
divisible by $p$ once $4k\ge3p$ (equality $4k=3p$ is impossible), so
exactly the range $m<k$, $4k<3p$ survives, giving \eqref{eq:Delta}.
The listed values are obtained by direct exact evaluation.
\end{proof}

\section{The Euler transfer modulo \texorpdfstring{$p^2$}{p-squared} and the pairing}\label{sec:4}

Write
\[
D_p(z)=H_p(z)-K_p(w(z))\in\Q[z],
\qquad w(z)=4z(1-z).
\]

\begin{theorem}[weighted transfer and first defect]\label{thm:C}
For every odd prime $p$,
\[
D_p(z)=
-\sum_{k=(p+1)/2}^{p-1}a_k64^{-k}
\sum_{j=p-k}^{k}(-1)^j\binom{k}{j}z^{k+j},
\]
hence $D_p\in p^2\Zp[z]$. Moreover, with $b=\fl{(3p-1)/4}$,
\begin{equation}\label{eq:defect}
\begin{gathered}
D_p(z)\equiv p^2z^pR_p(z)\pmod{p^3},\\
R_p(z)=-\sum_{k=(p+1)/2}^{b}\frac{a_k}{p^2}\,64^{-k}
\sum_{j=p-k}^{k}(-1)^j\binom{k}{j}z^{k+j-p},
\end{gathered}
\end{equation}
and, viewing the reductions $\bar R_p,\bar R_p'\in\Fp[z]$,
\begin{equation}\label{eq:refl}
\bar R_p(1-z)=-\bar R_p(z),\qquad \bar R_p'(1-z)=\bar R_p'(z)
\qquad\text{in }\Fp[z].
\end{equation}
\end{theorem}

\begin{proof}
By Theorem~\ref{thm:A}, the coefficients of degree $<p$ of $H-K\circ w$
vanish, so $D_p$ receives contributions only from terms of $K\circ w$
with $k>(p-1)/2$; expanding $(4z)^k(1-z)^k$ and discarding degrees
$<p$ gives the displayed formula (note $j+k\ge p$ is automatic there).
The integrality $D_p\in p^2\Zp[z]$ is immediate from
Lemma~\ref{lem:val}, including $p=3$. Reducing \eqref{eq:defect}: terms
with $v_p(a_k)\ge3$ (i.e. $k>b$) are $0$ modulo $p^3$, and the top
surviving exponent is $2b-p<p$, so $R_p$ has degree $2\fl{(p-3)/4}+1<p$.

For the reflection \eqref{eq:refl} we use the following elementary
projection lemma, applied repeatedly: a polynomial $T\in\Fp[z]$ of
degree $<2p$ decomposes \emph{uniquely} as
\begin{equation}\label{eq:proj}
T(z)=L(z)+z^pV(z),\qquad \deg L<p,\ \deg V<p,
\end{equation}
simply by splitting its monomials according to whether the exponent is
$<p$ or $\ge p$; if moreover $T(1-z)=T(z)$, then comparing the
coefficients of degree $\ge p$ in $T(1-z)=T(z)$ and using
$(1-z)^p=1-z^p$ in $\Fp[z]$ gives $V(1-z)=-V(z)$, so that
$T(1-z)-T(z)=-2z^pV(z)$ in $\Fp[z]$.

Apply this to the reduction $\bar D_p\in\Fp[z]$ of $D_p/p^2$, which has
degree $\le 2b-p+2(p-1)<2p$ by the degree bound above. Since
$w(1-z)=w(z)$, the polynomial $K_p\circ w$ is invariant under
$z\mapsto1-z$; in $\Fp[z]$ each term $z^k(1-z)^k$ of
$\bar D_p$ (coming from \eqref{eq:defect} modulo $p$, i.e. from
$z^k(1-z)^k$ expanded at degrees $\ge p$) is itself invariant, hence so
is $\bar D_p$; and $H_p$ is invariant only in the combination
$H_p(z)-H_p(1-z)$, so writing $\bar D_p=L+z^pV$ as in \eqref{eq:proj}
the invariance of each tail chunk $[z^k(1-z)^k]_{\ge p}=z^pV_k(z)$
gives $V_k(1-z)=-V_k(z)$ for every $k$, hence $V(1-z)=-V(z)$. But by
\eqref{eq:defect} the reduction of $V$ modulo $p$ is exactly
$\bar R_p$, the reduction of $R_p$. Therefore
$\bar R_p(1-z)=-\bar R_p(z)$ in $\Fp[z]$, which is the first identity
of \eqref{eq:refl}. Differentiating it gives the second. (Note that
\eqref{eq:refl} is genuinely a statement over $\Fp$: the first defect
$R_p$ enters $D_p$ only through the $p^2$-multiple in
\eqref{eq:defect}, so only its reduction modulo $p$ is defined
intrinsically.)
\end{proof}

\begin{corollary}[reflection-difference form]\label{cor:refldiff}
In $\Zp[z]$,
\begin{equation}\label{eq:refldiff}
H_p(z)-H_p(1-z)\equiv p^2R_p(z)\pmod{p^3}.
\end{equation}
\end{corollary}

\begin{proof}
Since $w(1-z)=w(z)$ we have
$H_p(z)-H_p(1-z)=D_p(z)-D_p(1-z)$ exactly. By \eqref{eq:defect},
\[
D_p(z)-D_p(1-z)\equiv
p^2\bigl(z^pR_p(z)-(1-z)^pR_p(1-z)\bigr)\pmod{p^3}.
\]
Reducing the right-hand factor modulo $p$ and using
$(1-z)^p\equiv1-z^p$ together with $R_p(1-z)\equiv-R_p(z)\pmod p$
from \eqref{eq:refl} gives
$z^pR_p(z)-(1-z)^pR_p(1-z)\equiv R_p(z)\pmod p$, and since the
congruence is $p^2$ times a quantity computed modulo $p$, this yields
\eqref{eq:refldiff}.
\end{proof}

\begin{corollary}[pairing]\label{cor:pair}
Let $\mathcal W_p(z)=zH_p'(z)=\sum_{n<p}nh_nz^n$. Then in $\Zp[z]$,
\begin{equation}\label{eq:pairing}
(1-z)\mathcal W_p(z)+z\mathcal W_p(1-z)
\equiv p^2z(1-z)R_p'(z)\pmod{p^3},
\end{equation}
and whenever $z,1-z$ are $p$-units,
\begin{equation}\label{eq:onesided}
\mathcal W_p(z)\equiv\frac{1-2z}{1-z}\,(\theta_wK_p)\bigl(w(z)\bigr)
\pmod{p^2},
\qquad \theta_w=w\frac{\dd}{\dd w}.
\end{equation}
\end{corollary}

\begin{proof}
Differentiating the congruence \eqref{eq:refldiff} of
Corollary~\ref{cor:refldiff} (differentiation is a well-defined
operation on coefficientwise congruences of polynomials) gives
$\mathcal W_p(z)/z+\mathcal W_p(1-z)/(1-z)\equiv p^2R_p'(z)\pmod{p^3}$;
multiplying by $z(1-z)$ yields \eqref{eq:pairing}. For
\eqref{eq:onesided}, differentiating $H_p\equiv K_p\circ w\pmod{p^2}$
and using $\theta_z w=4z(1-2z)$, $w/z=4(1-z)$ gives
$\mathcal W_p(z)\equiv\frac{1-2z}{1-z}(\theta_wK_p)(w(z))\pmod{p^2}$.
\end{proof}

\begin{remark}[sharpness]\label{rem:sharp}
The mod $p^2$ pairings of the table below do not lift to $\equiv0$
modulo $p^3$ in general: at $p=5$ one computes $R_5=z+2$, so
$R_5'\equiv1$, and the three denominator-cleared pairings
$\mathcal W_5(72)+8\mathcal W_5(576)$,
$\mathcal W_5(48)-4\mathcal W_5(-192)$,
$\mathcal W_5(63)-64\mathcal W_5(-4032)$ are congruent to
$50,75,75\pmod{125}$ respectively (direct evaluation), matching
Remark~\ref{rem:pairdefect}; all are nonzero. (Exact values in
Appendix~\ref{app:A}.)
\end{remark}

The three point pairs of Sun's conjecture arise as follows. Writing
\[
\mathcal W_p(M):=\sum_{n<p}n\binom{2n}{n}G_n\big/M^n=\mathcal W_p(64/M),
\]
and recalling $U_p(d)=\sum_{k<p}ka_k/d^k$ from \eqref{eq:TU}, one has
$w(z)/256=1/d$ and the Euler factors
$\frac{1-2z}{1-z}\in\{-7,\,5,\,65\}$:

\begin{center}
\footnotesize
\begin{tabular}{ccccc}
\toprule
$(M,64/(1-z))$ & $(z,1-z)$ & $d=256/w$ & Euler
& pairing (mod $p^2$)\\
\midrule
$(72,576)$ & $(\frac89,\frac19)$ & $648$ & $-7$ &
$\mathcal W_p(72)+8\mathcal W_p(576)\equiv0$\\
$(48,-192)$ & $(\frac43,-\frac13)$ & $-144$ & $5$ &
$\mathcal W_p(48)-4\mathcal W_p(-192)\equiv0$\\
$(63,-4032)$ & $(\frac{64}{63},-\frac1{63})$ & $-3969$ & $65$ &
$\mathcal W_p(63)-64\mathcal W_p(-4032)\equiv0$\\
\bottomrule
\end{tabular}
\end{center}

The first two pairings hold for $p>3$; the third requires also $p\ne7$
(both because $63$ must be a $p$-unit and because $1-z=-1/63$ appears).
In each case $1-z$ is a $p$-unit in the stated range, so
\eqref{eq:onesided} gives
\begin{equation}\label{eq:eulertrans}
\begin{gathered}
\mathcal W_p(72)\equiv-7U_p(648),\qquad
\mathcal W_p(48)\equiv5U_p(-144),\\
\mathcal W_p(63)\equiv65U_p(-3969)\pmod{p^2}.
\end{gathered}
\end{equation}

\begin{remark}[mod $p^3$ refinement of the pairings]\label{rem:pairdefect}
Clearing denominators in \eqref{eq:pairing} at the three points gives
the exact first defects of the paired sums:
\begin{align*}
\mathcal W_p(72)+8\mathcal W_p(576)
&\equiv \frac{8p^2}{9}\,R_p'\Bigl(\frac89\Bigr),\\
\mathcal W_p(48)-4\mathcal W_p(-192)
&\equiv \frac{4p^2}{3}\,R_p'\Bigl(\frac43\Bigr),\\
\mathcal W_p(63)-64\mathcal W_p(-4032)
&\equiv \frac{64p^2}{63}\,R_p'\Bigl(\frac{64}{63}\Bigr)
\end{align*}
modulo $p^3$ (in the ranges of the table, all displayed denominators
being $p$-units). At $p=5$, where $R_5=z+2$, the three right-hand
sides evaluate to $50,75,75\pmod{125}$ (cf.\ Remark~\ref{rem:sharp});
in particular the mod $p^2$ pairings are sharp.
\end{remark}

\section{Structure of the defect polynomial}\label{sec:5}

This section is not needed for the proof of the conjecture, but it
explains \emph{why} the first defect is manageable: $R_p$ satisfies a
constant-coefficient forcing equation and is carried by a kernel of
about half its apparent degree. Let $\varepsilon=\leg{-1}{p}$ and let
$\mathcal S$ be the operator \eqref{eq:S}.

\begin{theorem}\label{thm:D}
Let $p$ be an odd prime and $q=\fl{(p-3)/4}$.
\begin{enumerate}[label=\textup{(\roman*)},leftmargin=2.2em]
\item Forcing equation:
$\mathcal SR_p=\varepsilon/4$ in $\Fp[z]$.
\item Kernel: $R_p(z)=t\,J_p(w)$ with $t=1-2z$, $w=4z(1-z)$, where
\begin{equation}\label{eq:kernel}
J_p(w)=-\frac{\varepsilon}{2}\sum_{j=0}^{q}
\frac{(\frac12)_j(\frac34)_j(\frac54)_j}{(\frac32)_j^3}
\sum_{h=0}^{j}\frac{(-j)_h}{(\frac12-j)_h}w^h ,
\end{equation}
all denominators being $p$-units.
\item Integral invariant: with $a=\fl{p/4}$, $F=P_a(1-2z)$ the Legendre
polynomial part and $\lambda=a(a+1)$, the expression
\[
I_F[Y]=2s^2(F')^2Y-2s^2FF'Y'+s^2F^2Y''+stF^2Y'+2\lambda sF^2Y
\]
satisfies $I_F[Y]'=F^2\mathcal SY$ identically, and
\[
I_F[R_p]=\frac{\varepsilon}{4}\Bigl(\int_0^zF(u)^2\dd u-\frac{z^p}{2a+1}\Bigr)
\qquad\text{in }\Fp[z].
\]
The term $z^p/(2a+1)$ is forced by the fact that in characteristic $p$
the kernel of $\dd/\dd z$ on polynomials of degree $\le p$ is spanned by
$1$ and $z^p$; it is pinned down by $I_F[R_p](0)=I_F[R_p](1)=0$ and the
Legendre square integral $\int_0^1F^2=1/(2a+1)$.
\end{enumerate}
\end{theorem}

\begin{proof}
(i) Applying $\mathcal S$ to the truncations gives the exact boundary
computations
\[
\mathcal SH_p=-p^3h_pz^{p-1}+\frac{p(4p^2-1)}{4}h_{p-1}z^p,
\qquad
\mathcal MK_p=-p^3c_pw^{p-1},
\]
and the pullback \eqref{eq:pullback} holds for the truncations up to the
same boundary terms. Since $v_p(h_p),v_p(c_p)\ge1$, reducing
$\mathcal S(H_p-K_p\circ w)\equiv p^2\mathcal S(z^pR_p)\pmod{p^3}$
gives $\mathcal S(z^pR_p)=-(h_{p-1}/p)z^p/4$ in $\Fp[z]$. In
characteristic $p$ one has $\mathcal S(z^pY)=z^p\mathcal SY$ (every term
of $\mathcal S$ except the zeroth-order one contains a derivative), so
cancelling $z^p$ leaves $\mathcal SR_p=-(h_{p-1}/p)/4$. Finally
$h_{p-1}/p\equiv-\varepsilon\pmod p$: in the defining convolution for
$G_{p-1}$, every off-central summand contains a factor $p$ in at least
one of its central binomial coefficients.  The sole surviving term has
$k=m=(p-1)/2$; using
$\binom{p-1}{m}\equiv(-1)^m$ and $4^m\equiv(-1)^m$ shows that it equals
$(-1)^m=\varepsilon$ modulo $p$.  Thus
$G_{p-1}\equiv\varepsilon\pmod p$, and
$\binom{2p-2}{p-1}/p\equiv-1$ together with $64^{p-1}\equiv1$ yields
the sign.

 (ii) The reflection in \eqref{eq:refl} and $t^2=1-w$ first give a
unique factorization $R_p=tJ_p(w)$ with $\deg J_p\le q$.  To identify
$J_p$, put $m=(p-1)/2$, write $k=m+1+j$ and $r=p-1-k=m-1-j$, and set
\[
 V_{p,k}(z)=z^{-p}[z^k(1-z)^k]_{\ge p}.
\]
The coefficient of $z^{p-1}$ in the complementary low part is
$(-1)^r\binom kr\equiv\binom{2r}{r}=:B_r$.  Since the complete
polynomial $P=z^k(1-z)^k$ satisfies $sP'-ktP=0$, separating its low and
high parts gives the exact equation
\[
 sV_{p,k}'-ktV_{p,k}=-(2k+1)B_r.                 \tag{5.3}
\]
Writing $V_{p,k}=tQ_j(w)$ and using $k\equiv j+\frac12\pmod p$ turns
this into
\[
 w(1-w)Q_j'+(jw-j-\tfrac12)Q_j=-(2j+2)B_r.
\]
Its coefficient recursion, beginning with
$Q_j(0)=(2k+1)B_r/k$, has the polynomial solution
\[
 Q_j(w)=\frac{2k+1}{k}B_r
 \sum_{h=0}^j\frac{(-j)_h}{(\frac12-j)_h}w^h .  \tag{5.4}
\]
All denominators here are $p$-units.  The scalar of this tail in
\eqref{eq:defect} is
$u_j=(a_{m+1+j}/p^2)64^{-(m+1+j)}$.  Wilson's theorem gives
$u_0=1$ and
\[
 \frac{u_j}{u_{j-1}}=
 \frac{2j(16j^2-1)}{(2j+1)^3},\qquad
 \frac{((2k+1)/k)B_r}{((2k-1)/(k-1))B_{r+1}}
 =\frac{2j-1}{8j}.
\]
Their product is
$\frac{(j-1/2)(j-1/4)(j+1/4)}{(j+1/2)^3}$; the initial value is
$\varepsilon/2$.  Substitution in the defining tail for $R_p$ proves
\eqref{eq:kernel}.  Its top coefficient is nonzero, so
$\deg J_p=q$ and $\deg R_p=2q+1$.

(iii) The identity $I_F[Y]'=F^2\mathcal SY$ is a formal consequence of
$\mathcal DF=0$ via the product rules. Integrating (i) termwise
(legitimate since $2a+1<p$, so no denominator $p$ occurs) gives
\[
I_F[R_p]=\frac{\varepsilon}{4}\int_0^zF(u)^2\dd u+C_0+C_1z^p
\qquad\text{in }\Fp[z],
\]
where the integration ambiguity is $C_0+C_1z^p$ because in
characteristic $p$ the kernel of $\dd/\dd z$ on polynomials of degree
$\le p$ is spanned by $1$ and $z^p$. To pin down the constants we use
the reflections: $F(1-z)=(-1)^aF(z)$, and each of the five terms of
$I_F[Y]$ is built from an even number of $F$'s or $F'$'s together with
the reflection-odd $R_p$ (and $s=z(1-z)$ invariant, $t=1-2z$ odd), so
$I_F[R_p](1-z)=-I_F[R_p](z)$. Evaluation at $z=0$ gives
$I_F[R_p](0)=0$ (every term contains $s$ or $R_p(0)=0$), hence
$C_0=0$; evaluation at $z=1$ then gives
$0=I_F[R_p](1)=\frac{\varepsilon}{4}\int_0^1F^2+C_1$, and the
classical Legendre square integral (Rodrigues formula and integration
by parts, valid in $\Fp$ since $2a+1<p$)
$\int_0^1F(u)^2\dd u=1/(2a+1)$ yields $C_1=-\varepsilon/(4(2a+1))$.
\end{proof}

\begin{remark}
The forcing equation (i) with the reflection \eqref{eq:refl}
characterizes $R_p$ uniquely among polynomials of degree $<p$: the
homogeneous solution of $\mathcal SY=0$ regular at $0$ is $F^2$, which
is reflection-\emph{even}, while the recurrence for the coefficients is
determined by the constant term. Formula \eqref{eq:kernel} and the
invariant (iii) are two cross-checkable presentations of the same
defect.
\end{remark}

\section{Special values: the three weighted pairs}\label{sec:6}

We now combine \eqref{eq:eulertrans} with published special values of
the triple binomial sums \eqref{eq:TU}. The small primes excluded by the
printed hypotheses are settled by direct exact computation
(Appendix~\ref{app:A}).

\subsection{The pair \texorpdfstring{$(72,576)$}{(72,576)}}

\begin{theorem}\label{thm:grpA}
For every prime $p>3$,
\[
\mathcal W_p(72)\equiv-8\mathcal W_p(576)\equiv
\begin{cases}
4x^2-3p,&p=x^2+4y^2\equiv1\pod4,\\
p,&p\equiv3\pod4,
\end{cases}
\pmod{p^2}.
\]
\end{theorem}

\begin{proof}
The pairing is Corollary~\ref{cor:pair}. For $p\ne7$, Sun
\cite[Theorem~2.10]{SunBQF} gives (the printed hypotheses exclude
$p=2,3,7$)
\[
U_p(648)\equiv
\begin{cases}
(3p-4x^2)/7\pmod{p^2},&p=x^2+4y^2\equiv1\pmod4,\\
-p/7\pmod{p^2},&p\equiv3\pmod4,
\end{cases}
\]
and \eqref{eq:eulertrans} with factor $-7$ yields the values. For
$p=7$: direct evaluation from \eqref{eq:Gdef} gives
$\mathcal W_7(72)\equiv7\pmod{49}$ and
$\mathcal W_7(72)+8\mathcal W_7(576)\equiv0\pmod{49}$
(Appendix~\ref{app:A}), which is the $p\equiv3\pmod4$ branch.
\end{proof}

\subsection{The pair \texorpdfstring{$(48,-192)$}{(48,-192)}}

\begin{theorem}\label{thm:grpB}
For every prime $p>3$,
\[
\mathcal W_p(48)\equiv4\mathcal W_p(-192)\equiv
\begin{cases}
3p-4x^2,&p=x^2+3y^2\equiv1\pod3,\\
-p,&p\equiv2\pod3,
\end{cases}
\pmod{p^2}.
\]
\end{theorem}

\begin{proof}
For $p>5$, Sun \cite[Theorem~2.9]{SunBQF} gives $U_p(-144)$ modulo
$p^2$ with values $(3p-4x^2)/5$ and $-p/5$ in the split and inert cases;
multiplication by $5$ via \eqref{eq:eulertrans} gives the claim.
The printed range of that theorem is $p>3$, but the formula containing
$1/5$ does not hold at $p=5$ as stated: $U_5(-144)\equiv19\pmod{25}$
while $-p/5\equiv-1\equiv24\pmod{25}$. We therefore quote it for $p>5$
only and settle $p=5$ directly: $\mathcal W_5(48)\equiv20\equiv-5
\pmod{25}$ and $\mathcal W_5(48)-4\mathcal W_5(-192)\equiv0\pmod{25}$
(Appendix~\ref{app:A}), the inert branch.
\end{proof}

\subsection{The pair \texorpdfstring{$(63,-4032)$}{(63,-4032)}: a CM Ramanujan sum}

By \eqref{eq:eulertrans}, the third pair reduces to $65U_p(-3969)$.
The required input is the following explicit specialization of the
general CM theorem of Chisholm--Deines--Long--Nebe--Swisher
\cite[Theorem~1]{CDLNS} (restated with both reduction types in
\cite[Theorem~2.1]{BRSTT}).

\begin{proposition}\label{prop:ram65}
For every prime $p>3$, $p\ne7$,
\begin{equation}\label{eq:ram65}
\sum_{k=0}^{p-1}(65k+8)\frac{a_k}{(-3969)^k}
\equiv8\leg{-7}{p}p\pmod{p^2}.
\end{equation}
\end{proposition}

\begin{proof}
\emph{Parameters.} Since
$\frac{(1/2)_k(1/4)_k(3/4)_k}{(k!)^3}=\frac{a_k}{256^k}$, the sum
\eqref{eq:ram65} is $8$ times the normalized truncated sum with
\[
d=4,\qquad \lambda=-\frac{256}{3969},\qquad \alpha=\frac{65}{8},
\qquad \delta=\frac{9\sqrt7}{8};
\]
the identity $\sum_{k\ge0}(65k+8)a_k/(-3969)^k=9\sqrt7/\pi$
\cite[Table~2]{Guillera} certifies that these are the parameters of a
Ramanujan-type formula in the sense required by the cited theorem. (An
independent normalization check: Clausen--Pfaff gives
$K(-4x/(1-x)^2)=\sqrt{1-x}\,F(x)$ between the ${}_3F_2$ with parameters
$(\frac14,\frac12,\frac34)$ and $(\frac12,\frac12,\frac12)$, and at
$x=1/64$ differentiation gives
$(65\theta_w+8)K=\frac{9\sqrt7}{16}(42\theta_x+5)F$; the classical
value $(42k+5)$-series $=16/\pi$ \cite[(1.1)]{OZ} gives the same
$\delta$.)

\emph{Hypotheses.} $1-\lambda=(65/63)^2$, so $\Q(\sqrt{1-\lambda})=\Q$
(totally real, unramified conditions automatic) and $|\lambda|<1$.
With $t=(1-\sqrt{1-\lambda})/2=-1/63$, the curve of the theorem is
\[
E:\ y^2=x\Bigl(x^2+x-\tfrac1{252}\Bigr),\qquad
\Delta(E)=\frac{2^6}{3^6 7^3},\qquad j(E)=255^3=16581375.
\]
The $j$-invariant is the singular modulus of the CM order of
discriminant $-28$ (the class polynomial is $H_{-28}(X)=X-16581375$;
compare \cite{LMFDB}), so $E$ has CM by $\Q(\sqrt{-7})$; for $p>3$,
$p\ne7$ the displayed model has $p$-integral coefficients and unit
discriminant, hence good reduction. By Deuring's criterion the ordinary
(resp. supersingular) primes are the split (resp. inert) primes of
$\Q(\sqrt{-7})$, so the sign in the theorem is
$\operatorname{sgn}(E,p)=\leg{-7}{p}$. For $p\notin\{2,3,5,7,13\}$ both
$\alpha$ and $\lambda$ are $p$-units and $\leg{1-\lambda}{p}=1$, so
\cite[Theorem~1]{CDLNS} applies verbatim and multiplication by $8$
gives \eqref{eq:ram65} --- including the inert classes
$p\equiv3,5,6\pmod7$, for which the theorem assigns $-p$ in the
normalized sum (we do not use any ordinary-prime-only variant).

\emph{The two exceptional primes.} At $p=5,13$ the parameter
$\alpha=65/8$ is not a $p$-unit and $p^2\mid1-\lambda$, so the theorem
cannot be invoked. Direct exact evaluation of the cleared-denominator
sum (no division by $65$ is ever performed) gives
\[
\begin{array}{c|c|c}
p & \bigl((65k+8)a_k(-3969)^{-k}\bmod p^2\bigr)_{0\le k<p} & \text{sum}\\
\hline
5 & 8,\,17,\,10,\,0,\,0 & 10\equiv-40\pmod{25}\\
13 & 8,\,160,\,142,\,80,\,52,\,65,\,65,\,0,\dots,0 & 65\equiv-104\pmod{169}
\end{array}
\]
Both primes are inert in $\Q(\sqrt{-7})$, and the sums equal
$-8p$ modulo $p^2$, as required. (Terms with $k>(p-1)/2$ vanish modulo
$p^2$ by Lemma~\ref{lem:val}.)
\end{proof}

\begin{remark}[a period-product form]\label{rem:period}
Proposition~\ref{prop:fclausen} allows formal differentiation modulo
$p^2$: at $t_0=-1/63$ the chain rule reads $65\theta_w=64\theta_t$,
and with $B_p$ as in Section~\ref{sec:2},
\[
\sum_{k<p}(65k+8)\frac{a_k}{(-3969)^k}
\equiv8B_p(t_0)\bigl(B_p(t_0)+16\,\theta_tB_p(t_0)\bigr)\pmod{p^2}
\]
for $p\nmid2\cdot3\cdot5\cdot7\cdot13$; the two factors are the finite
period and its Gauss--Manin derivative of $E$. This gives a concrete
geometric reading of \eqref{eq:ram65}; it is not used in the proof.
\end{remark}

\begin{theorem}\label{thm:grpC}
For every prime $p>3$, $p\ne7$,
\[
\mathcal W_p(63)\equiv64\mathcal W_p(-4032)\pmod{p^2},
\]
with common value
\[
\begin{cases}
-8(4x^2-3p),&p=x^2+7y^2\equiv1,2,4\pod7,\\
-8p,&p\equiv3,5,6\pod7.
\end{cases}
\]
\end{theorem}

\begin{proof}
Write $T_p=T_p(-3969)$, $U_p=U_p(-3969)$.
Proposition~\ref{prop:ram65} reads $65U_p+8T_p\equiv8\leg{-7}{p}p$.
For split $p$ ($p=x^2+7y^2\equiv1,2,4\pmod7$), Sun--Ye
\cite[Theorem~1.8]{SunYe} give $T_p\equiv4x^2-2p\pmod{p^2}$ (in fact
modulo $p^3$); for inert $p$ ($p\equiv3,5,6\pmod7$, $p\ne7$), Sun
\cite[Theorem~5.1]{SunL2} gives $T_p\equiv0\pmod{p^2}$. Hence
\[
65U_p\equiv8\leg{-7}{p}p-8T_p\equiv
\begin{cases}
8p-8(4x^2-2p)=-8(4x^2-3p),&\text{split},\\
-8p,&\text{inert},
\end{cases}
\]
and \eqref{eq:eulertrans} identifies the left side with
$\mathcal W_p(63)\equiv64\mathcal W_p(-4032)\pmod{p^2}$.
\end{proof}

\section{The denominator 128 modulo \texorpdfstring{$p^3$}{p-cubed}}\label{sec:7}

\begin{theorem}[split branch]\label{thm:128split}
For $p=x^2+2y^2\equiv1,3\pmod8$,
\[
\sum_{n=0}^{p-1}\frac{\binom{2n}{n}G_n}{128^n}
\equiv4x^2-2p-\frac{p^2}{4x^2}\pmod{p^3}.
\]
\end{theorem}

\begin{proof}
Theorem~\ref{thm:B} reduces the left side to
$T_p=\sum_{k<p}a_k/256^k$ modulo $p^3$, and Sun--Ye
\cite[Theorem~1.5]{SunYe} evaluate exactly this sum:
$T_p\equiv4x^2-2p-p^2/(4x^2)\pmod{p^3}$ for all
$p=x^2+2y^2\equiv1,3\pmod8$ (no excluded primes). Note $p=3=1^2+2\cdot1^2$
is covered; the central binomial $\binom{2n}{n}G_n/128^n$ is
$p$-integral throughout.
\end{proof}

For the inert classes we use the jet comparison proved in
Appendix~\ref{app:B} and the following exact rational identity, proved
in Appendix~\ref{app:C}. Let, for $a\ge1$,
\[
f_i=(-1)^i\binom ai\binom{a+i}{i},
\]
\[
\mathcal Q_a(z)=\sum_{j=0}^{a-1}\ \sum_{k=0}^{a-1-j}
f_{k+j+1}\,(-1)^k\frac{(a-k-1)!\,(k!)^2}{(a+k+1)!}\,z^j .
\]

\begin{theorem}\label{thm:E}
For every odd $a\ge1$,
\[
\mathcal Q_a\Bigl(\tfrac12\Bigr)
=-\Bigl(\frac{2^{a-1}}{a\binom{a-1}{(a-1)/2}}\Bigr)^2 .
\]
\end{theorem}

\begin{proposition}[jet reduction at the center]\label{prop:jet}
Let $p\ge5$ and $a=\fl{p/4}$. Then
$H_p(\tfrac12)\equiv A_p(\tfrac12)^2+\frac{3p^2}{16}\mathcal
Q_a(\tfrac12)\pmod{p^3}$,
where $A_p$ is the truncation of $B$. If $a$ is odd
(equivalently $p\equiv5,7\pmod8$), then $A_p(\tfrac12)\in p^2\Zp$, so
\[
\sum_{n=0}^{p-1}\frac{\binom{2n}{n}G_n}{128^n}
\equiv\frac{3p^2}{16}\mathcal Q_a\Bigl(\tfrac12\Bigr)\pmod{p^3}.
\]
\end{proposition}

\begin{theorem}[inert branches]\label{thm:128inert}
Let $p>3$ be prime and $C_p=\binom{\lfloor p/4\rfloor}{\lfloor p/8\rfloor}$.
Then
\[
\sum_{n=0}^{p-1}\frac{\binom{2n}{n}G_n}{128^n}
\equiv
\begin{cases}
\dfrac{p^2}{3}C_p^{-2}\pmod{p^3},&p\equiv5\pmod8,\\[2mm]
-\dfrac{3p^2}{2}C_p^{-2}\pmod{p^3},&p\equiv7\pmod8.
\end{cases}
\]
\end{theorem}

\begin{proof}
Write $p=8k+5$ or $8k+7$; then $a=\fl{p/4}=2k+1$, $\fl{p/8}=k$, and
$C_p=\binom{a}{k}$, a $p$-unit. Theorem~\ref{thm:E} gives
$\mathcal Q_a(\frac12)=-16^k/((k+1)^2\binom{a}{k}^2)$.
If $p=8k+5$: $16^k=2^{(p-5)/2}=2^{(p-1)/2}/4\equiv\leg2p/4=-\frac14$
and $k+1\equiv\frac38\pmod p$; if $p=8k+7$: $16^k\equiv\frac18$ and
$k+1\equiv\frac18$. Substituting into Proposition~\ref{prop:jet}:
\[
\frac{3}{16}\mathcal Q_a\Bigl(\tfrac12\Bigr)\equiv
\begin{cases}
\frac{3}{16}\cdot\frac{1/4}{(3/8)^2}C_p^{-2}=\frac13C_p^{-2},&p\equiv5\pmod8,\\[1mm]
\frac{3}{16}\cdot\frac{-1/8}{(1/8)^2}C_p^{-2}=-\frac32C_p^{-2},&p\equiv7\pmod8,
\end{cases}
\]
and multiplication by $p^2$ gives the claim.
\end{proof}

\section{Conclusion and remarks}\label{sec:8}

\begin{theorem}\label{thm:final}
Conjecture~2.4 of \cite{Sun24} holds: the three weighted pairs are
Theorems~\ref{thm:grpA}, \ref{thm:grpB}, \ref{thm:grpC}, and the
denominator-$128$ congruences are Theorems~\ref{thm:128split} and
\ref{thm:128inert}.
\end{theorem}

\begin{remark}[higher precision at the center]
Theorem~\ref{thm:Bp} shows that the center bridge lifts to $p^4$
exactly at the primes where the finite sum \eqref{eq:Delta} vanishes
modulo $p$; the values $1,1,2,1,5,12$ at $p=3,5,7,11,13,17$ show this
is not automatic. We have no criterion for the vanishing in general.
\end{remark}

\begin{remark}[no naive finite Pfaff transplant]\label{rem:pfaff}
One might hope to convert \eqref{eq:ram65} into the classical $42k+5$
truncated sum via a direct finite transformation. This is impossible
without a correction term: the known modulo-$p^2$ value
$\sum_{k<p}(42k+5)\binom{2k}{k}^3/4096^k\equiv5\leg{-1}{p}p$
(compare \cite[Theorem~1.1]{OZ}) is incompatible with
\eqref{eq:ram65} under the infinite-series prefactor $9\sqrt7/16$
unless $p\mid47^2$. Any finite analogue of the Clausen--Pfaff
transformation between these two truncations must therefore contain a
first-order tail correction; we have not constructed one.
\end{remark}

\begin{remark}[scope]
The same two-level transfer applies in principle to other CM points of
the family, provided the corresponding special values of the triple
binomial sums are known; the defect analysis of Section~\ref{sec:5} is
point-free. We have not pursued further instances.
\end{remark}

\appendix

\section{Exceptional primes: exact certificates}\label{app:A}

All computations below are exact integer/rational evaluations of the
defining sums \eqref{eq:Gdef}, \eqref{eq:TU}; they are finite
certificates covering exactly the primes excluded by the quoted
theorems, not samples suggesting a general pattern.

\medskip
\noindent\textbf{Group $(72,576)$ at $p=7$.}
$\mathcal W_7(72)\equiv7\pmod{49}$ and
$\mathcal W_7(576)\equiv0\pmod{49}$; in particular the pairing and the
value $p$ hold.

\noindent\textbf{Group $(48,-192)$ at $p=5$.}
$\mathcal W_5(48)\equiv20\equiv-5\pmod{25}$ and
$\mathcal W_5(-192)\equiv0\pmod{25}$; the pairing and the value $-p$
hold. For comparison, the printed $p>3$ formula of
\cite[Theorem~2.9]{SunBQF} would give $U_5(-144)\equiv-1\pmod{25}$,
while the exact value is $19$; the corrected combination
$5U_5(-144)+T_5(-144)$ is still $\equiv-5\pmod{25}$.

\noindent\textbf{Proposition~\ref{prop:ram65} at $p=5,13$.} The full
residue tables are displayed in the proof of
Proposition~\ref{prop:ram65}; the sums are $10\equiv-40\pmod{25}$ and
$65\equiv-104\pmod{169}$, matching $8\leg{-7}{p}p$ since both primes
are inert in $\Q(\sqrt{-7})$.

\noindent\textbf{Center at $p=3$.}
$\sum_{n<3}\binom{2n}{n}G_n/128^n\equiv\sum_{k<3}a_k/256^k\equiv16
\pmod{27}$; the exact difference is $945/8192=3^3\cdot35/8192$ with
$\Delta_3=35/8192\equiv1\pmod3$, matching \eqref{eq:Delta}.

\noindent\textbf{Pairing defects at $p=5$} (cf. Remark~\ref{rem:sharp}):
\begin{align*}
\mathcal W_5(72)+8\mathcal W_5(576)&\equiv50,\\
\mathcal W_5(48)-4\mathcal W_5(-192)&\equiv75,\\
\mathcal W_5(63)-64\mathcal W_5(-4032)&\equiv75
\end{align*}
modulo $125$.

\section{The parametric jet and the corrected reflection}\label{app:B}

This appendix proves Proposition~\ref{prop:jet}. Fix $p\ge5$,
$a=\fl{p/4}$, $b=p-a-1$. Let $B_n=\binom{2n}{n}\binom{4n}{2n}$,
$A_p(z)=\sum_{n<p}B_nz^n/64^n$, and define the finite Pochhammer
polynomial
\[
J_p(\nu,z)=\sum_{n=0}^{p-1}\frac{(-\nu)_n(\nu+1)_n}{(n!)^2}z^n
\in\Q(\nu)[z],
\]
whose Taylor coefficients at $\nu=a$ are $p$-integral since each $n!$ is
a $p$-unit. Write
$J_p(a+s,z)=F(z)+sU(z)+s^2V(z)+O(s^3)$.

\begin{lemma}\label{lem:B1}
$F(z)=P_a(1-2z)$ has degree $a$, and
$A_p\equiv F-\frac p4U+\frac{p^2}{16}V\pmod{p^3}$ in $\Zp[z]$.
\end{lemma}

\begin{proof}
At $\nu=a$,
\[
 \frac{(-\nu)_n(\nu+1)_n}{(n!)^2}
 =(-1)^n\frac{(a+n)!}{(a-n)!(n!)^2}qquad(0\le n\le a),
\]
which is the coefficient of $P_a(1-2z)$; for $n>a$ the factor
$(-\nu)_n$ vanishes.  Thus $F=P_a(1-2z)$ exactly.  Now
\[
 a-\frac p4=
 \begin{cases}-\frac14,&p\equiv1\pmod4,\\
 -\frac34,&p\equiv3\pmod4,
 \end{cases}
\]
and in the two cases the upper parameters
$(-(a-p/4),a-p/4+1)$ are respectively
$(\frac14,\frac34)$ and $(\frac34,\frac14)$.  Consequently
\[
 A_p(z)=J_p(a-p/4,z)
\]
coefficientwise, with no analytic continuation involved.  Taylor's
formula at the integral point $a$ therefore gives
\[
 A_p=F-\frac p4U+\frac{p^2}{16}V+p^3Z
\]
for some $Z\in\Zp[z]$: every coefficient of $J_p$ is a polynomial in
$\nu$ over $\Zp$, because $n!$ is a $p$-unit for $n<p$.  This proves
the asserted congruence and also fixes the signs and the normalization
of $V$ (the coefficient of $s^2$, not the second derivative itself).
\end{proof}

\begin{lemma}\label{lem:B2}
$[z^n]U$ is given, with $f_n=(-1)^n\binom an\binom{a+n}{n}$, by
\[
[z^n]U=
\begin{cases}
f_n(H_{a+n}-H_{a-n}),&0\le n\le a,\\
(-1)^{a+1}(n-a-1)!\,(n+a)!/(n!)^2,&a<n<p,
\end{cases}
\]
and $D_p:=U\bmod p$ has degree $b$. With
$W_p(z)=\sum_{n=b+1}^{p-1}\bigl([z^n]U/p\bigr)z^n\bmod p$, Wilson's
theorem rewrites it as the short reversed polynomial
\[
W_p(z)=\sum_{k=0}^{a-1}w_kz^{p-1-k},
\qquad
w_k=(-1)^k\frac{(a-k-1)!(k!)^2}{(a+k+1)!}.
\]
\end{lemma}

\begin{proof}
For $n\le a$, logarithmic differentiation of both Pochhammer factors
at $\nu=a$ gives
\[
 \sum_{r=0}^{n-1}\frac1{a-r}
 +\sum_{r=1}^{n}\frac1{a+r}=H_{a+n}-H_{a-n},
\]
which proves the first line.  For $n>a$, only the factor
$a-\nu$ in $(-\nu)_n$ vanishes at $a$; differentiating it and
multiplying the remaining factors gives the second line, including the
sign $(-1)^{a+1}$.

If $n>b=p-a-1$, the factor $(\nu+1)_n$ contains exactly one multiple of
$p$ at $\nu=a$, while all its other factors are $p$-units.  Locally the
two distinguished factors are $(-s)(p+s)$.  If the product of all
remaining factors is $K(s)\in\Zp[s]$, then
$[z^n]U=-pK(0)$ and $[z^n]V=-K(0)-pK'(0)$; hence
$[z^n]V\equiv[z^n]U/p\pmod p$.  In particular these coefficients of
$U$ are divisible by $p$, so $D_p=U\bmod p$ has degree at most $b$.

Finally put $n=p-1-k$, $0\le k<a$.  Wilson's theorem gives
\[
 \begin{aligned}
 \frac{(n+a)!}{p}&\equiv-(a-k-1)!,&
 (n-a-1)!&\equiv\frac{(-1)^{a+k}}{(a+k+1)!},\\
 (n!)^2&\equiv\frac1{(k!)^2}&&\pmod p.
 \end{aligned}
\]
Combining these congruences with the factor $(-1)^{a+1}$ in the high
coefficient of $U$ yields precisely the coefficient $w_k$ displayed
above.  This proves the reversed expression for $W_p$.
\end{proof}

\begin{lemma}[error decomposition]\label{lem:B3}
With $[\,\cdot\,]_{\ge p}$ denoting truncation to degrees $\ge p$ and
$E_p=(H_p-A_p^2)/p^2\bmod p$,
\[
E_p=\Bigl[\tfrac38FW_p-\tfrac1{16}D_p^2\Bigr]_{\ge p}\pmod p .
\]
\end{lemma}

\begin{proof}
By Theorem~\ref{thm:C} and Proposition~\ref{prop:fclausen},
$H_p-A_p^2\in p^2\Zp[z]$, so $E_p$ is a well-defined polynomial in
$\Fp[z]$, supported in degrees $\ge p$ (the coefficients of degree $<p$
of $H_p-A_p^2$ vanish exactly by Theorem~\ref{thm:A}). Square the
expansion of Lemma~\ref{lem:B1}: the $-\frac p2FU$ term contributes
$-[FW]/2$ after truncation (the low part of $U$ does not overflow since
$\deg F+\deg D=a+b=p-1$), and the $\frac{p^2}{16}(U^2+2FV)$ term
contributes $[D^2]/16+[FW]/8$ using $[z^n]V\equiv[z^n]U/p$ on the
overflow range; $E_p$ being the \emph{negative} of the overflow of
$A_p^2$ gives the displayed coefficients $3/8$ and $-1/16$.
\end{proof}

\begin{lemma}[reflection of the derivative part]\label{lem:B4}
$D_p(1-z)=(-1)^aD_p(z)$ in $\Fp[z]$.
\end{lemma}

\begin{proof}
With $\mathcal L_a=z(1-z)\partial_z^2+(1-2z)\partial_z+a(a+1)$, the
coefficient recursion of the finite Pochhammer polynomial, together with
the $p$-divisibility of the high coefficients of $U$, gives
$\mathcal L_aD_p=-(2a+1)F$ in $\Fp[z]$; the last coefficient equation
is consistent since $\deg D_p=b\le p-2$. Among polynomials of degree
$<p$, $\ker\mathcal L_a$ is spanned by $F$ (the recursion
$(n+1)^2Y_{n+1}+(a(a+1)-n(n+1))Y_n=0$ determines $Y$ from $Y_0$, all
$(n+1)^2$ being nonzero). Since $\mathcal L_a$ commutes with
$\rho:z\mapsto1-z$ and $\rho F=(-1)^aF$ (Rodrigues), the difference
$\rho D_p-(-1)^aD_p$ is both a multiple of $F$ and anti-invariant under
$\rho$ up to the sign $(-1)^a$; in odd characteristic this forces it to
vanish.
\end{proof}

\begin{proof}[Proof of Proposition~\ref{prop:jet}]
Apply the projection lemma to $D_p^2$: a polynomial
$T\in\Fp[z]$ of degree $<2p$ with $T(1-z)=T(z)$ writes uniquely as
$L+z^pC$ with $\deg L,\deg C<p$, and comparing degrees $\ge p$ via
$(1-z)^p=1-z^p$ gives $C(1-z)=-C(z)$. Combined with Lemma~\ref{lem:B4}
and the reversed form of $W_p$, the product $[FW_p]_{\ge p}$ is exactly
$z^p\mathcal Q_a(z)$ with $\mathcal Q_a$ as in Section~\ref{sec:7}
(the entries of $\mathcal Q_a$ are the convolution of the $f_i$ and
$w_k$), while $[D_p^2]_{\ge p}$ has antisymmetric $z^p$-part. Hence
\[
\widetilde R_p:=E_p z^{-p}-\tfrac38\mathcal Q_a
\quad\text{satisfies}\quad
\widetilde R_p(1-z)=-\widetilde R_p(z),
\]
and evaluating at $z=\frac12$ (note $(\frac12)^p=\frac12$) gives
$E_p(\frac12)=\frac3{16}\mathcal Q_a(\frac12)$; since
$H_p-A_p^2=p^2E_p$ in $\Zp[z]$,
$H_p(\frac12)\equiv A_p(\frac12)^2+\frac{3p^2}{16}\mathcal
Q_a(\frac12)\pmod{p^3}$.
If $a$ is odd, $F(\frac12)=P_a(0)=0$ exactly and
$D_p(\frac12)=0$ by Lemma~\ref{lem:B4}, so
$A_p(\frac12)\equiv\frac{p^2}{16}V(\frac12)\pmod{p^3}$ lies in
$p^2\Zp$ and its square is $0$ modulo $p^3$.
\end{proof}

\section{Proof of Theorem~\ref{thm:E}: the symmetric-square
identity}\label{app:C}

Fix $a\ge1$ odd, $\lambda=a(a+1)$, $s=z(1-z)$, $t=1-2z$, and let
$F=F_a(z)=\sum_{i=0}^a f_iz^i$ with
$f_i=(-1)^i\binom ai\binom{a+i}{i}$, so $\mathcal DF=0$ for
$\mathcal D=s\partial_z^2+t\partial_z+\lambda$, and let
\[
\mathcal S=s^2\partial^3+3st\,\partial^2
+\bigl(t^2+(4\lambda-2)s\bigr)\partial+2\lambda t
\]
be the symmetric-square operator, satisfying \eqref{eq:symmsq} with
respect to $\mathcal D$.

\emph{Step 1: $\mathcal Q:=\mathcal Q_a$ as a positive part.}
Set $w_{j-1}=-1/(j^2f_j)$ (consistent with Appendix~\ref{app:B}) and
\[
V(z)=\sum_{j=1}^{a}\frac{z^{-j}}{j^2f_j},
\qquad\text{so}\qquad
\mathcal Q=-[FV]_+,
\]
where $[\,\cdot\,]_+$ takes the nonnegative-degree part. The recursion
$(j+1)^2f_{j+1}=-(a-j)(a+j+1)f_j$ gives the exact Laurent identity
\begin{equation}\label{eq:DV}
\mathcal DV=-z^{-1}+f_a^{-1}z^{-a-1}.
\end{equation}

\emph{Step 2: forcing term.} Substituting \eqref{eq:DV} into
\eqref{eq:symmsq} with $f=F$, $v=V$: the term $-z^{-1}$ contributes
$F-3(1-z)F'$, and the nonnegative part of the $f_a^{-1}z^{-a-1}$
contribution is the constant $-(2a+1)$ (read off the top coefficients
of $3sF'$, $tF$, $(a+1)(1-z)F$: $-3a,-2,a+1$). The operator $\mathcal S$
preserves the purely negative part (its possible top-degree coefficient
on $z^n$ is $(n+1)(n-2a)(n+2a+2)$, vanishing at $n=-1$ and negative
beyond), so the positive-part extraction is compatible and
\begin{equation}\label{eq:SQ}
\mathcal S\mathcal Q=3(1-z)F'-F+(2a+1).
\end{equation}

\emph{Step 3: symmetrization.} With $q=(\mathcal Q+\rho\mathcal Q)/2$,
$\rho Y(z)=Y(1-z)$: for odd $a$, $\rho F=-F$ (Rodrigues) and
$\mathcal S\rho=-\rho\mathcal S$; taking the symmetric part of
\eqref{eq:SQ} kills the constant and gives
\begin{equation}\label{eq:Sq}
\mathcal Sq=\tfrac32tF'-F .
\end{equation}

\emph{Step 4: the invariant.} For any polynomial $Y$, the expression
\[
J[Y]=2s^2(F')^2Y-2s^2FF'Y'+s^2F^2Y''+stF^2Y'+2\lambda sF^2Y
\]
satisfies $J[Y]'=F^2\mathcal SY$ identically (a formal consequence of
$\mathcal DF=0$: with $F,F',F''$ treated as independent jet variables,
the difference of the two sides factors exactly as
$(4sF'Y-2sFY')(sF''+tF'+\lambda F)$). Hence by \eqref{eq:Sq},
\[
(2J[q])'=3tF^2F'-2F^3=(tF^3)',
\]
so $2J[q]=tF^3+C$; at $z=0$ one has $s=0$, $F=1$, $t=1$, and every term
of $J[q]$ contains $s$, so $C=-1$:
\begin{equation}\label{eq:Jq}
2J[q]=tF^3-1 .
\end{equation}

\emph{Step 5: evaluation.} At $z=\frac12$: $s=\frac14$, $t=0$, and
$F(\frac12)=P_a(0)=0$ ($a$ odd), so \eqref{eq:Jq} reduces to
$4s^2(F')^2q=-1$; since $q(\frac12)=\mathcal Q(\frac12)$,
\[
\mathcal Q_a\Bigl(\tfrac12\Bigr)=-\frac{4}{F_a'(1/2)^2}.
\]
Finally $F_a'(1/2)=-2P_a'(0)$. The coefficient of $x$ in the Legendre
polynomial $P_{2m+1}(x)$ is
\[
[x]\,P_{2m+1}(x)=(-1)^m(2m+1)\binom{2m}{m}\big/4^m;
\]
thus, with $a=2m+1$,
\[
\mathcal Q_a\Bigl(\tfrac12\Bigr)
=-\frac{16^m}{(2m+1)^2\binom{2m}{m}^2}
=-\Bigl(\frac{2^{a-1}}{a\binom{a-1}{(a-1)/2}}\Bigr)^2 .
\qed
\]

\section{Computational checks}\label{app:D}

All symbolic claims in this paper are proved in full generality in the
text. In addition, the following were verified by exact machine
computation (Python, standard-library rational arithmetic, from the
defining convolution \eqref{eq:Gdef}; independent PARI/GP
reimplementations agree): the coefficients of Theorem~\ref{thm:A} up to
degree $24$; the identity \eqref{eq:pascal} for $r\le24$; the center
bridge and the defect formula \eqref{eq:Delta} for the $24$ odd primes
$p\le97$; the transfer, reflection, forcing equation, kernel and
integral invariant of Sections~\ref{sec:4}--\ref{sec:5} for
$3\le p\le31$; the values in Appendix~\ref{app:A}; and all nine
formulas of the conjecture for the primes $5\le p\le43$ including the
quadratic-form reconstructions. These checks are diagnostic; no general
statement in the paper rests on them.

Selected algebraic components of the argument (the operator pullback
\eqref{eq:pullback} after clearing denominators, the dyadic recursion
behind \eqref{eq:pascal}, and the divisibility assemblies of
Sections~\ref{sec:6}--\ref{sec:7}) have been checked in Lean; the proof
as a whole has not been formalized.

\bibliographystyle{amsplain}
\bibliography{references}

\end{document}
